% Method - ICML style, why-driven narrative

\subsection{Factoring Route Generation into Commitment and Execution}

Route generation under corridor multi-modality requires resolving two fundamentally different types of uncertainty: \emph{which corridor} to take (a discrete, global choice) and \emph{how to traverse} the chosen corridor (a continuous, local process).
CascadeTraj makes this factorization explicit:
\begin{equation}
  p_\theta(\mathbf{w} \mid c) = \int p_\theta(z \mid c) \cdot p_\phi(\mathbf{w} \mid z, c) \, dz,
  \label{eq:factorization}
\end{equation}
where $z \in \mathbb{R}^D$ is a corridor-level latent commitment sampled in the \emph{decision stage} $p_\theta(z \mid c)$, and $\mathbf{w}$ is the way-level route produced in the \emph{execution stage} $p_\phi(\mathbf{w} \mid z, c)$.
The decision stage operates in a low-dimensional space shaped to represent corridor identity; the execution stage operates on the road graph, constrained to produce topologically valid paths.

\subsection{Way-Level Tokenization: The Right Abstraction}
\label{sec:way-tokenization}

The choice of route representation determines both the sequence length the model must handle and the structure of the output space.
Node-level tokenization produces sequences of hundreds of steps, where exposure bias compounds over the full route and the output vocabulary spans tens of thousands of nodes.
GPS-coordinate representations avoid the discrete vocabulary problem but sacrifice road-network structure, producing trajectories that may drift off-road.

Way-level tokenization resolves both issues.
An OSM way is a maximal road segment of uniform type and name---the semantic atom of road network topology.
Tokenizing a route as a way sequence naturally compresses $L \sim$ hundreds of nodes to $M \sim$ tens of ways, because each way typically spans multiple consecutive nodes.
Crucially, adjacency is preserved: two consecutive ways in a valid route share at least one intersection node, so the space of valid next tokens at each step is determined by the road graph rather than by the full vocabulary.

\subsection{Decision Stage: Corridor Commitment via Latent Flow}

\subsubsection{Learning the Corridor Space}

The autoencoder maps variable-length way sequences to a fixed latent representation and back.

\paragraph{Way encoder.}
Each way $w_m$ is embedded as a $d_{\mathrm{model}}$-dimensional token using learned embeddings for way identity, road tier, and positional features (grid coordinates, length).
A Perceiver encoder \citep{jaegle2021perceiver} compresses the variable-length way token sequence into $n_L = 8$ fixed latent tokens via cross-attention, where 8 learnable query tokens attend to all way tokens.
The number 8 reflects a prior that a typical urban route contains roughly 8 major turning points or corridor transitions; the Perceiver's cross-attention adaptively extracts the most salient information from routes of any length.

\paragraph{Information bottleneck.}
The $n_L \times d_{\mathrm{model}} = 2048$-dimensional Perceiver output cannot be used directly as the flow model's target space.
In corridor multi-modal settings, the latent representations of different corridors occupy isolated regions of high-dimensional space with topologically infeasible interpolations between them---unlike images, where interpolating two latent codes produces a plausible blend, interpolating two corridor codes yields a ``ghost route'' on no valid road.
A flow model trained in this 2048-dimensional space tends toward mode averaging, converging to the conditional mean $\mathbb{E}[z \mid c]$ that falls in the void between corridors.

A $\beta$-VAE bottleneck addresses this by projecting the Perceiver output into a $D$-dimensional space:
\begin{equation}
  z_{\text{flat}} \in \mathbb{R}^{2048}
  \xrightarrow{\mathrm{MLP}}
  (\boldsymbol{\mu}, \log \boldsymbol{\sigma}^2) \in \mathbb{R}^D \times \mathbb{R}^D
  \xrightarrow{\text{reparam}}
  z_{\text{vec}} \in \mathbb{R}^D
  \xrightarrow{\mathrm{MLP}}
  \hat{z}_{\text{flat}} \in \mathbb{R}^{2048}.
  \label{eq:bottleneck}
\end{equation}
The training loss combines reconstruction and regularization:
\begin{equation}
  \mathcal{L}_{\mathrm{AE}} = \mathcal{L}_{\mathrm{CE}}(\hat{\mathbf{w}}, \mathbf{w}) + \beta \cdot D_{\mathrm{KL}}(q(z \mid \mathbf{w}) \| p(z)),
  \label{eq:ae-loss}
\end{equation}
with $\beta$ warmed up from 0 to 0.01 over 30 epochs.
The bottleneck serves two purposes: (i)~it forces the latent to retain corridor identity while discarding intra-corridor variation, because the reconstruction loss preserves what matters and the limited capacity discards what does not; (ii)~the KL regularization shapes the posterior mean $\boldsymbol{\mu}$ toward a near-Gaussian distribution, making it a tractable target for flow-based generation.

\paragraph{Capacity matching.}
The bottleneck dimension $D$ must match the intrinsic dimensionality of the corridor space.
At $D\!=\!64$, the KL divergence stabilizes around 6~nats ($\approx$8.7~bits), implying an effective state count of $\sim\!400$---consistent with the diversity of corridor alternatives in a city-scale road network.
Section~\ref{sec:ablation-dim} shows that deviations from this matched capacity produce asymmetric degradation: $D\!=\!32$ causes mild precision loss, while $D\!=\!128$ triggers catastrophic mode averaging.

\subsubsection{Conditional Flow Matching on $\boldsymbol{\mu}$}
\label{sec:flow}

The flow model generates corridor commitments by learning a conditional velocity field that transforms a standard Gaussian into the distribution of posterior means $\boldsymbol{\mu} \in \mathbb{R}^{64}$.
We use $\boldsymbol{\mu}$ rather than the sampled $z_{\text{vec}}$ because $\boldsymbol{\mu}$ is the deterministic summary of the posterior---free of sampling noise and more stable as a regression target \citep{lipman2023flow}.

Trip context $c = (o, d, t)$ is injected through cross-attention layers in a 6-layer transformer.
At inference, $K$ latent samples $\{\boldsymbol{\mu}_k\}_{k=1}^K$ are drawn by integrating the learned velocity field from Gaussian noise, then mapped back to $n_L \times d_{\mathrm{model}}$ latent tokens via the bottleneck decoder MLP.

\subsection{Execution Stage: Constrained Autoregressive Decoding}
\label{sec:decoder}

Each latent sample $\boldsymbol{\mu}_k$ specifies a corridor commitment; the decoder's task is to realize it as a graph-feasible way sequence.

\paragraph{Graph-constrained candidate scoring.}
Rather than predicting among all $\sim$35K ways, the decoder scores only the graph neighbors of the current way position---at most 32 candidates at each step.
This constraint guarantees graph feasibility by construction and compresses the effective output space by three orders of magnitude, substantially reducing exposure bias.

\paragraph{Corridor guidance via cross-attention.}
The decoder attends to the $n_L\!=\!8$ latent tokens through cross-attention at each decoding step, receiving corridor-level directional guidance that resolves ambiguity among local candidates.
A candidate-aware query mechanism lets the decoder jointly consider all neighbor candidates in context before scoring.

\paragraph{Anti-loop decoding.}
Corridor commitment resolves global direction, but local decoding can still produce short-range loops (oscillating between two or three adjacent ways).
At inference time, a log-penalty of $\alpha\!=\!2.0$ is applied to any way that appeared in the most recent $k\!=\!4$ steps.
This correction requires no training modification and no additional parameters; it improves arrival rate from 70.2\% to 78.3\% (Section~\ref{sec:ablation-antiloop}).

\subsection{Inference Pipeline}
\label{sec:inference}

Given conditioning $c$, inference proceeds in three steps:
(1)~the flow model samples $K\!=\!16$ corridor commitments $\{\boldsymbol{\mu}_k\}$;
(2)~the decoder generates a candidate route for each $\boldsymbol{\mu}_k$ using graph-constrained autoregressive decoding with anti-loop;
(3)~among successful candidates (those reaching~$d$), the shortest route is selected for deployment, or all $K$ routes are retained for coverage and diversity evaluation.
